\documentclass{scrartcl}
\usepackage{graphicx}  % required for inserting images

% for cyrillic and Bulgarian language
\usepackage[utf8]{inputenc}
\usepackage[T2A]{fontenc}
\usepackage[bulgarian]{babel}

% math packages
\usepackage{amsmath}
\usepackage{float}
\usepackage{amsfonts}
\usepackage{amssymb}
\usepackage{amsthm}
\usepackage{cancel}

\usepackage{ragged2e}  % adds FlushLeft

\usepackage{listings}
\usepackage{xcolor}

\definecolor{codegreen}{rgb}{0,0.6,0}
\definecolor{codegray}{rgb}{0.5,0.5,0.5}
\definecolor{codepurple}{rgb}{0.58,0,0.82}
\definecolor{backcolour}{rgb}{0.95,0.95,0.92}

\lstdefinestyle{mystyle}{
    label={lst:listing-cpp},
    language={C++},
    backgroundcolor=\color{backcolour},   
    commentstyle=\color{codegreen},
    keywordstyle=\color{magenta},
    numberstyle=\tiny\color{codegray},
    stringstyle=\color{codepurple},
    basicstyle=\ttfamily\footnotesize,
    breakatwhitespace=false,         
    breaklines=true,                 
    keepspaces=true,                  
    numbers=left,                    
    numbersep=5pt,                  
    showspaces=false,                
    showstringspaces=false,
    showtabs=false,                  
    tabsize=2,
}
\lstset{style=mystyle}

\usepackage[colorlinks=true, linkcolor=blue, urlcolor=cyan]{hyperref}

\newtheorem{theorem}{Теорема}
\newtheorem{definition}{Дефиниция}

\title{Компютърни архитектури}
\author{Ивайло Андреев + gemini 3 Flash}
\date{31 май 2026}

\begin{document}
\maketitle

\tableofcontents
\newpage

\begin{FlushLeft}

\section{Обща структура на компютрите и концептуално изпълнение на инструкциите, запомнена програма} 

\subsection{Обща структура на компютрите} 
Персоналният компютър (ПК) представлява универсална цифрова изчислителна система, предназначена за решаване на разнообразни алгоритмични задачи и проблеми. Основните му специфични особености включват компактни размери, ограничена и проста система от команди, ограничен обем на основната памет и опростен интерфейс за свързване на компонентите.

Въпреки голямото разнообразие от производители, персоналните компютри се характеризират с еднакъв модел на вътрешната архитектура, базиран на класическия модел на Фон Нойман, който се състои от три основни хардуерни компонента:
\begin{itemize}
    \item \textbf{Централен процесор (ЦП / CPU)}: Представлява основната изчислителна сила на компютъра, която изпълнява всички аритметични и логически операции и управлява жизнения цикъл на инструкциите.
    \item \textbf{Оперативна памет (ОП / RAM)}: Енергозависима памет, която позволява бързо четене и писане. В нея се съхраняват изпълняваните в момента програми и файлове с данни, с които ЦП работи директно. При архитектурата на Фон Нойман както данните, така и програмите се съхраняват в едно и също споделено пространство на ОП.
    \item \textbf{Входно-изходни (периферни) устройства (ВИ)}: Служат за предаване или получаване на информация от/към външната среда (клавиатура, мишка, монитор, мрежова карта NIC и др.).
\end{itemize}

\subsubsection{Организация на комуникацията: Интерфейсна шина} 
Обменът на сигнали, команди и данни между компонентите може да се организира по два начина:
\begin{enumerate}
    \item \textbf{Индивидуални информационни канали (връзка "всеки с всеки")}: Компонентите образуват пълен граф на свързаност. Скоростта е висока, но сложността и цената нарастват неимоверно с добавянето на нови устройства.
    \item \textbf{Обща информационна шина}: Споделена комуникационна магистрала, достъпна за всички устройства в конкурентен режим. Тя е по-евтина и лесна за надграждане, но в даден момент могат да комуникират само две устройства, което може да се превърне в теснолинейка на производителността.
\end{enumerate}

Поради конкурентния достъп се въвежда специализирано хардуерно устройство – \textbf{арбитър (контролер) на шината}, който управлява приоритетите и определя кое устройство кога да заеме линията. На ЦП се дава най-висок приоритет. 

Логически шината се разделя на три подшини според функциите си:
\begin{itemize}
    \item \textbf{Подшина за данни}: Пренася информационните битове между ЦП, паметта и ВИ. Тя е двупосочна. Ширината ѝ определя броя битове, пренасяни едновременно.
    \item \textbf{Подшина за адреси}: Еднопосочна подшина, определяща конкретното физическо местоположение в паметта или ВИ. Нейната ширина определя максималното количество адресируема памет (напр. 16-битова шина адресира $2^{16} = 64$ KB, а 32-битова – 4 GB).
    \item \textbf{Управляваща подшина}: Пренася електрически контролни сигнали (за четене/писане, линии за синхронизация, заявки за прекъсвания и сигнали за статус).
\end{itemize}

\subsection{Концептуално изпълнение на инструкциите} 
Изпълнението на всяка машинна инструкция в процесора преминава последователно през следните пет базови етапа:
\begin{enumerate}
    \item \textbf{Извличане (Fetch)}: Инструкцията се прочита от ОП от адреса, указан в програмния брояч ($IP / PC$).
    \item \textbf{Дешифриране (Decode)}: Кодът на операцията се анализира от управляващия блок и понякога се разбива на елементарни микрооперации ($\mu$-ops).
    \item \textbf{Осигуряване на операнди (Fetch Operands)}: Изчисляват се адресите на входните данни и те се зареждат от паметта или регистрите.
    \item \textbf{Изпълнение (Execute)}: Аритметико-логическото устройство извършва съответната операция.
    \item \textbf{Запис на резултата (Writeback)}: Получената стойност се записва в целевия регистър или в ОП.
\end{enumerate}

\subsubsection{Оптимизационни хардуерни стратегии}
Съвременните суперскаларни процесори използват паралелизъм на ниво инструкции за ускоряване на работата:
\begin{itemize}
    \item \textbf{Конвейерна обработка (Pipelining)}: Припокриване на етапите на изпълнение. Докато една инструкция се изпълнява, следващата се дешифрира, а по-следващата се извлича. Това предотвратява престоя на хардуерните блокове.
    \item \textbf{Изпълнение извън реда (Out-of-Order Execution)}: Инструкциите се пренареждат динамично така, че тези, които имат готови операнди, да се изпълнят веднага, без да чакат забавени предходни инструкции.
    \item \textbf{Предсказване на преходи (Branch Prediction)}: Статистическо отгатване на посоката на условните преходи с цел конвейерът да не спира работа при чакане на оценка на флаговете.
    \item \textbf{Спекулативно изпълнение (Speculative Execution)}: Изпълнение на инструкции по предположения път преди преходът да е реално изчислен. При грешка резултатите се отхвърлят.
    \item \textbf{Преименуване на регистри (Register Renaming)}: Използване на допълнителни скрити физически регистри за елиминиране на фалшиви зависимости между данните (архитектурни конфликти).
\end{itemize}

\subsection{Запомнена програма} 
Концепцията за запомнена програма изисква кодът да се съхранява в паметта като поредица от числа. При \textbf{архитектурата на Фон Нойман} програмите се зареждат от постоянни носители (ROM, HDD/SSD) в общата ОП преди изпълнение. При \textbf{Харвард архитектурата} съществуват две физически разделени памети – една изрично за инструкции и една за данни.

В съвременните системи всяка програма съдържа една или повече \textbf{нишки (threads)}, които се изпълняват конкурентно и могат да бъдат прекъсвани (от ОС или хардуерни събития). За да бъде това прекъсване напълно "прозрачно" за програмата, хардуерът и ОС съхраняват т.нар. \textbf{векторно състояние на процеса} (съдържание на програмния брояч, стековия указател, регистрите с данни и управляващите таблици). След обработка на прекъсването векторът се възстановява и процесът продължава без загуба на контекст. Важно правило в софтуерното инженерство е, че изпълнимият код не трябва да се видоизменя динамично по време на работа (самомодифициращ се код се счита за лош стил).

\section{Формати на данните} 

В паметта на компютъра всички данни се представят в двоичен вид (като битове 0 и 1). Чрез $n$-битово поле могат да се кодират точно $2^n$ различни състояния или стойности.

\subsection{Цели двоични числа} 
Целите числа се представят в машинни думи или байтове с фиксирана дължина (8, 16, 32, 64 бита). Те биват два вида:
\begin{itemize}
    \item \textbf{Без знак (Unsigned)}: Всички $n$ бита се използват за представяне на абсолютната стойност. Диапазонът е от $0$ до $2^n - 1$.
    \item \textbf{Със знак (Signed)}: Най-старшият бит ($a_{n-1}$) се отделя за знак (0 за положително, 1 за отрицателно число).
\end{itemize}

За кодиране на отрицателни числа исторически съществуват три основни метода:
\begin{enumerate}
    \item \textbf{Прав код (Sign-Magnitude)}: Старшият бит е знак, останалите са модулът на числото. Диапазонът е $[-2^{n-1}+1, 2^{n-1}-1]$. Недостатък: Нулата има две кодирания ($+0$ и $-0$), а аритметичните операции се реализират много трудно апаратно.
    \item \textbf{Обратен код (Ones' Complement)}: Отрицателните числа се образуват чрез инвертиране (обръщане на 0 в 1 и обратно) на правия код. Нулата отново има две представяния, а при пренос извън най-старшия бит при събиране, към крайния резултат трябва циклично да се добави 1.
    \item \textbf{Допълнителен код / Допълнение към двойката (Two's Complement)}: Най-широко използваната схема в съвременните процесори. Отрицателно число се получава чрез инвертиране на битовете на положителното и добавяне на единици ($+1$) към най-младшия бит. 
\end{enumerate}

Математически, стойността на едно $N$-битово число в допълнителен код се изчислява по следната формула:
\[ \overline{a_{N-1} a_{N-2} \dots a_0}_{2c} = -a_{N-1} \cdot 2^{N-1} + \sum_{i=0}^{N-2} a_i \cdot 2^i \]
Следователно диапазонът е асиметричен: $[-2^{N-1}, 2^{N-1} - 1]$. Нулата има само едно единствено уникално представяне ($\overline{00\dots0}_{2c}$). Предимството на допълнителния код е, че хардуерът извършва събиране и изваждане по един и същ начин, а преносът извън най-старшия бит просто се игнорира, което улеснява логическите схеми.

\subsection{Двоично-десетични числа (BCD)} 
При BCD формата (\textit{Binary-Coded Decimal}) всяка десетична цифра от 0 до 9 се представя самостоятелно чрез 4-битово двоично поле (тетрада). Тъй като с 4 бита могат да се представят 16 комбинации, 6 от тях остават неизползвани (невалидни). Съществуват две форми за организация в байтове:
\begin{itemize}
    \item \textbf{Непакетиран BCD формат}: В един байт се съхранява само една десетична цифра (обикновено в по-младата тетрада, а старшата се запълва с нули).
    \item \textbf{Пакетиран BCD формат}: В един байт се съхраняват две десетични цифри (една в старшата и една в младшата тетрада).
\end{itemize}
\textbf{Предимство}: Липсват хардуерни загуби от закръгляне при конвертиране на десетични дробни числа в двоични (напр. числото $0.3$ е безкрайна периодична двоична дроб, но в BCD се кодира точно). \textbf{Недостатък}: Форматът е силно неефективен по отношение на паметта и изчисленията, изисква допълнителни коригиращи инструкции след аритметични операции.

\subsection{Двоични числа с плаваща запетая} 
Рационалните реални числа се представят в компютъра чрез експоненциален запис с плаваща запетая (\textit{Floating-Point}). Те апроксимират реалните числа и се състоят от три компонента: знак на числото ($s$), мантиса ($m$) и експонента ($e$). Стойността се определя като:
\[ \text{Value} = (-1)^s \times m \times 2^e \] 

Съвременните архитектури имплементират международния стандарт \textbf{IEEE 754}, който дефинира следните основни формати:
\begin{itemize}
    \item \textbf{binary32 (Single Precision / float)}: Заема общо 32 бита – 1 бит за знак, 8 бита за експонента и 23 бита за мантиса.
    \item \textbf{binary64 (Double Precision / double)}: Заема общо 64 бита – 1 бит за знак, 11 бита за експонента и 52 бита за мантиса.
\end{itemize}
Мантисата е нормализирана (във формат $1.f$), като водещата единица пред десетичната запетая не се пази физически в паметта (скрит бит), което добавя допълнителна точност. Експонентата се съхранява с отместване (bias), за да се избегне поддържането на знак в нея.

\subsection{Символни данни и кодови таблици} 
Символите и текстовете се представят като поредици от числа посредством стандартизирани \textbf{кодови таблици}:
\begin{itemize}
    \item \textbf{ASCII}: 7-битов (или разширен 8-битов) стандарт. Кодовете от 0 до 127 са твърдо дефинирани за латиницата, цифрите и контролните символи. Регионът $[128, 255]$ зависи от локалната кодова страница (напр. CP1251 за кирилица), което често води до проблеми с нечетими текстове при грешна интерпретация.
    \item \textbf{Unicode (UTF-8, UTF-16, UTF-32)}: Универсална система, обхващаща всички писмености. \textbf{UTF-8} е с променлива дължина (от 1 до 4 байта) и е напълно обратно съвместим с ASCII за първите 127 символа. \textbf{UTF-16} и \textbf{UTF-32} използват фиксирани или по-големи базови структури от 2 или 4 байта.
\end{itemize}

\section{Вътрешна структура на централен процесор} 

Централният процесор съдържа вътрешни подблокове, които си взаимодействат по вътрешнопроцесорна шина:

\subsection{Регистри} 
Регистрите представляват свръхбързи паметови клетки, разположени директно в ядрото на процесора. Техният размер варира от 8 до 512 бита, като стандартно съвпадат с разрядността на архитектурата (машинната дума). Класифицират се като:
\begin{itemize}
    \item \textbf{Регистри с общо предназначение (GPR)}: Използват се от програмиста за съхранение на междинни данни при операции (в x86 това са EAX, EBX, ECX, EDX и техните 64-битови аналози RAX, RBX...).
    \item \textbf{Програмен брояч / Указател на инструкции ($IP / PC$)}: Пази адреса на следващата инструкция, която трябва да бъде извлечена от паметта.
    \item \textbf{Стеков указател ($SP / ESP$)}: Сочи към върха на системния стек в ОП.
    \item \textbf{Сегментни и контролни регистри}: Управляват режимите на работа на процесора и адресацията на паметта (CS, DS, SS).
\end{itemize}

\subsection{Аритметико-логическо устройство (АЛУ)} 
АЛУ е комбинационна логическа схема, която извършва същинската обработка на данните. То приема операнди от вътрешните регистри и изпълнява над тях аритметични (събиране, изваждане, умножение) и логически операции (AND, OR, NOT, XOR, битови измествания).

\subsection{Регистър на състоянието и флагове (Flags Register)} 
Това е специализиран регистър, в който всеки отделен бит действа като самостоятелен двупозиционен превключвател (\textbf{флаг}). Флаговете отразяват резултата от последното действие, извършено в АЛУ, и служат като вход за условните преходи. Основни флагове са:
\begin{itemize}
    \item \textbf{ZF (Zero Flag)}: Вдига се в 1, ако резултатът от операцията е нула.
    \item \textbf{CF (Carry Flag)}: Указва пренос или заем от най-старшия бит при операции без знак.
    \item \textbf{OF (Overflow Flag)}: Указва аритметично препълване при операции с допълнителен код (числа със знак).
    \item \textbf{SF (Sign Flag)}: Равен е на най-старшия бит на резултата (индикира отрицателно число).
\end{itemize}

\subsection{Блок за управление (Управляващо устройство - УУ)} 
Блокът за управление е "мозъкът" на процесора. Той извлича инструкциите, дешифрира ги и генерира поредица от вътрешни микрокоманди (управляващи импулси), които синхронизират работата на АЛУ, регистрите и шинните интерфейси в съответствие с тактовия генератор.

\section{Инструкции на централен процесор} 

Машинната инструкция е завършена команда за процесора, представена в двоичен вид.

\subsection{Структура на машинната инструкция}
Една типична инструкция (например в x86 архитектурата) се състои от следните логически полета:
\begin{itemize}
    \item \textbf{Префикси}: Незадължителни байтове, поставени преди кода на операцията, които модифицират поведението ѝ (например префикси за повторение при низове, пренасочване на сегменти или промяна на размера на операнда).
    \item \textbf{Код на операцията (Opcode)}: Задължително поле, което идентифицира конкретната команда (например събиране, преход, преместване).
    \item \textbf{Спецификатори на операндите}: Полета, указващи местоположението на входните и изходните данни (регистър, адрес в паметта или директна стойност).
\end{itemize}

\subsection{Модели на адресация на операндите}
Адресацията дефинира начина, по който се изчислява реалният (ефективният) адрес на данните в паметта:
\begin{itemize}
    \item \textbf{Директна (непосредствена) адресация}: Операндът е част от самата инструкция (напр. \texttt{MOV EAX, 5} – числото 5 е записано директно в кода).
    \item \textbf{Регистърна адресация}: Данните се намират в посочен регистър (\texttt{MOV EAX, EBX}).
    \item \textbf{Пряка (абсолютна) адресация}: Инструкцията съдържа точния числен адрес от ОП (\texttt{MOV EAX, [1000H]}).
    \item \textbf{Косвена регистърна адресация}: Регистър съдържа адреса към клетката от паметта (\texttt{MOV EAX, [EBX]}).
    \item \textbf{Базова индексирана адресация с отместване}: Адресът се изчислява динамично като сума от базов регистър, индексен регистър (евентуално мащабиран) и константно отместване (напр. при работа с масиви: \texttt{MOV EAX, [EBX + ESI*4 + 10H]}).
\end{itemize}

\subsection{Функционални групи инструкции}

\subsubsection{Аритметико-логически инструкции}
Извършват обработка в АЛУ. Примери в x86:
\begin{itemize}
    \item \texttt{ADD / SUB}: Събиране и изваждане на операнди.
    \item \texttt{MUL / IMUL}: Умножение без знак и със знак (изискват детайлно внимание към разрядността и знаковия бит).
    \item \texttt{AND / OR / XOR / NOT}: Побитови логически операции.
    \item \texttt{SHL / SHR / SAR}: Логически и аритметични измествания на битове наляво/надясно.
\end{itemize}

\subsubsection{Низови инструкции (String Instructions)}
Специализирани команди за обработка на последователности от данни в паметта (байтове, думи). Те автоматично инкрементират или декрементират индексните регистри (ESI и EDI) в зависимост от флага за посока (DF):
\begin{itemize}
    \item \texttt{MOVS}: Копира елемент от низа, сочен от ESI, на адреса в EDI.
    \item \texttt{CMPS}: Сравнява елементи от два низа.
    \item \texttt{SCAS}: Сравнява елемент от низ със съдържанието на акумулатора (AL/AX/EAX).
    \item \texttt{LODS / STOS}: Зарежда в / записва от акумулатора символ в низа.
\end{itemize}
Тези инструкции често се комбинират с префикси за повторение (\texttt{REP, REPE, REPNE}), които зациклят изпълнението им, докато броячът ECX стане равен на 0 или се наруши условието за равенство.

\subsubsection{Безусловни и условни преходи} 
Променят стандартния последователен поток на програмата чрез модифициране на програмния брояч ($IP / EIP$):
\begin{itemize}
    \item \textbf{Безусловен преход (\texttt{JMP})}: Винаги прехвърля изпълнението към посочения адрес.
    \item \textbf{Условни преходи (\texttt{JE, JNE, JG, JL} и др.)}: Изпълняват се само ако е изпълнено конкретно логическо условие, проверено чрез текущото състояние на флаговия регистър (напр. \texttt{JE} скача, ако $ZF = 1$).
\end{itemize}

\subsubsection{Управление на програмата}
Инструкции, които организират жизнения цикъл на програмната структура, функциите и прекъсванията:
\begin{itemize}
    \item \texttt{CALL}: Извиква подпрограма (функция). Преди прехода записва текущия адрес за връщане в системния стек.
    \item \texttt{RET}: Връща изпълнението от подпрограмата, като извлича адреса за връщане от стека.
    \item \texttt{INT / INTO}: Предизвиква софтуерно прекъсване (преход към системна подпрограма на ОС).
    \item \texttt{HLT}: Спира изпълнението на инструкции от процесора и го привежда в чакащо състояние до следващ външен хардуерен импулс (прекъсване).
\end{itemize}

\end{FlushLeft}

\end{document}